\documentclass[10pt,a4paper]{amsart}
\usepackage{amsmath,amssymb,amsthm,graphicx}
% Retained preamble for the Stacks Project chapter {CH} (shared project preamble).
% The chapter file pinned-original-source.tex only carries \input{{preamble}} on its first line,
% so no preamble text lives inside the chapter itself.  The Stacks Project uses one single shared
% preamble for every chapter: the snapshot stages it three times (sources/stacks-chow/preamble.tex,
% sources/stacks-etale-cohomology/preamble.tex and focused-seeds/AGP/sources/stacks-snapshot-20261011/
% preamble.tex) and those copies are identical in size.  The block below is copied mechanically,
% line for line, from sources/stacks-chow/preamble.tex lines 14-18 and 154-217: the Xy-pic setup and
% the native theorem declarations together with the mathematical macros.  Nothing was added, removed
% or redefined; the wrapper class and the amsmath/amssymb/amsthm/graphicx base packages are supplied
% by the assembler.
% For commutative diagrams we use Xy-pic
% Chapter: stacks-cohomology.tex -- Cohomology of Sheaves.
\usepackage[all]{xy}

% We use 2cell for 2-commutative diagrams.
\xyoption{2cell}

% Theorem environments.
%
\theoremstyle{plain}
\newtheorem{theorem}[subsection]{Theorem}
\newtheorem{proposition}[subsection]{Proposition}
\newtheorem{lemma}[subsection]{Lemma}

\theoremstyle{definition}
\newtheorem{definition}[subsection]{Definition}
\newtheorem{example}[subsection]{Example}
\newtheorem{exercise}[subsection]{Exercise}
\newtheorem{situation}[subsection]{Situation}

\theoremstyle{remark}
\newtheorem{remark}[subsection]{Remark}
\newtheorem{remarks}[subsection]{Remarks}

\numberwithin{equation}{subsection}

% Macros
%
\def\lim{\mathop{\mathrm{lim}}\nolimits}
\def\colim{\mathop{\mathrm{colim}}\nolimits}
\def\Spec{\mathop{\mathrm{Spec}}}
\def\Hom{\mathop{\mathrm{Hom}}\nolimits}
\def\Ext{\mathop{\mathrm{Ext}}\nolimits}
\def\SheafHom{\mathop{\mathcal{H}\!\mathit{om}}\nolimits}
\def\SheafExt{\mathop{\mathcal{E}\!\mathit{xt}}\nolimits}
\def\Sch{\mathit{Sch}}
\def\Mor{\mathop{\mathrm{Mor}}\nolimits}
\def\Ob{\mathop{\mathrm{Ob}}\nolimits}
\def\Sh{\mathop{\mathit{Sh}}\nolimits}
\def\NL{\mathop{N\!L}\nolimits}
\def\CH{\mathop{\mathrm{CH}}\nolimits}
\def\proetale{{pro\text{-}\acute{e}tale}}
\def\etale{{\acute{e}tale}}
\def\QCoh{\mathit{QCoh}}
\def\Ker{\mathop{\mathrm{Ker}}}
\def\Im{\mathop{\mathrm{Im}}}
\def\Coker{\mathop{\mathrm{Coker}}}
\def\Coim{\mathop{\mathrm{Coim}}}

% Boxtimes
%
\DeclareMathSymbol{\boxtimes}{\mathbin}{AMSa}{"02}

%
% Macros for moduli stacks/spaces
%
\def\QCohstack{\mathcal{QC}\!\mathit{oh}}
\def\Cohstack{\mathcal{C}\!\mathit{oh}}
\def\Spacesstack{\mathcal{S}\!\mathit{paces}}
\def\Quotfunctor{\mathrm{Quot}}
\def\Hilbfunctor{\mathrm{Hilb}}
\def\Curvesstack{\mathcal{C}\!\mathit{urves}}
\def\Polarizedstack{\mathcal{P}\!\mathit{olarized}}
\def\Complexesstack{\mathcal{C}\!\mathit{omplexes}}
% \Pic is the operator that assigns to X its picard group, usage \Pic(X)
% \Picardstack_{X/B} denotes the Picard stack of X over B
% \Picardfunctor_{X/B} denotes the Picard functor of X over B
\def\Pic{\mathop{\mathrm{Pic}}\nolimits}
\def\Picardstack{\mathcal{P}\!\mathit{ic}}
\def\Picardfunctor{\mathrm{Pic}}
\def\Deformationcategory{\mathcal{D}\!\mathit{ef}}
\usepackage{cleveref}
\crefname{originalreference}{original reference}{original references}
\Crefname{originalreference}{Original reference}{Original references}
\makeatletter
\newlabel{more-algebra-item-shift-tensor}{{0FNH}{source}{}{sourcepointer-1}{}}
\newlabel{more-algebra-item-shift-tensor@cref}{{[originalreference][0][]0FNH}{[1][0][]source}{}{}{}}
\makeatother
\title{The adjunction-constructed product on hypercohomology is the cup product}
\author{The Stacks Project Authors}
\date{stacks-project:cohomology.tex:local-snapshot-20261011; GNU FDL 1.2 or later}
\begin{document}
\maketitle
\noindent\textit{Editorial note.} All delivered mathematical text is the author's own, copied from the pinned Stacks Project chapter file cohomology.tex; no statement or proof step was rewritten, reordered or removed.\par\medskip
\noindent\textit{Editorial note.} The statement is a single sentence referring to the construction in the paragraph immediately before the lemma; that paragraph lies outside the counted span, and the construction is also written out inside the proof, so nothing load bearing is left outside.\par\medskip
\noindent\textit{Editorial note.} The author's footnote isolating the sign hidden in the shift-tensor identification is part of the counted proof span and is reproduced as written.\par\medskip
\noindent\textit{Editorial note.} The counted spans carry 1 navigation stub and no bibliography citation.\par\medskip
\noindent\textit{Editorial note.} The preamble is the chapter's retained minimal preamble (sources/stacks-cohomology/retained-preamble.tex lines 1-81), a verbatim excerpt of the native preamble.\par\medskip
\noindent\textit{Editorial note.} The label lemma-second-cup-equals-first is defined only in cohomology.tex, so no cross-chapter restatement had to be weighed for this card.\par\medskip
\noindent\textit{Allowed standard prerequisites (editorial).}\begin{itemize}
\item Adjunction between derived pullback and derived pushforward, and its counit. Derived Categories, the adjunction Lf\textasciicircum{}* and Rf\_* with counit epsilon, used at proof lines 7222-7229. Supplies the description of the constructed product as the adjoint of a composite of counits.
\item Shift-tensor identification for complexes of modules. More on Algebra, Item more-algebra-item-shift-tensor, invoked at proof line 7240. Supplies the identification A[-i] tensor A[-j] = A[-i-j] used to compare the two maps.
\item Derived global sections as the derived pushforward to a point. Same chapter, the identification Rf\_*(-) = RGamma(X, -) for f: X to a point, used at proof lines 7220-7221. Supplies the setting in which the cup product is defined.
\end{itemize}\medskip
\ifdefined\hypertarget\hypertarget{sourcepointer-1}{}\fi
\noindent\textit{Source reference.} Original reference 0FNH: sources/stacks-chow/tags.txt, source lines 19567--19567. Pointer outside the counted span, invoked once for context; the load-bearing argument of the counted statement is the author's proof reproduced in full in the proof span.\par\medskip
\begin{lemma}
\label{lemma-second-cup-equals-first}
This construction gives the cup product.
\end{lemma}

\begin{proof}
With $f : (X, \mathcal{O}_X) \to (pt, A)$ as above we have
$Rf_*(-) = R\Gamma(X, -)$ and our map $\mu$ is adjoint to the map
$$
Lf^*(Rf_*K \otimes_A^\mathbf{L} Rf_*M) =
Lf^*Rf_*K \otimes_{\mathcal{O}_X}^\mathbf{L} Lf^*Rf_*M
\xrightarrow{\epsilon_K \otimes \epsilon_M}
K \otimes_{\mathcal{O}_X}^\mathbf{L} M
$$
where $\epsilon$ is the counit of the adjunction between
$Lf^*$ and $Rf_*$.
If we think of $\xi$ and $\eta$ as maps $\xi : A[-i] \to R\Gamma(X, K)$
and $\eta : A[-j] \to R\Gamma(X, M)$, then
the tensor $\xi \otimes \eta$ corresponds to the map\footnote{There
is a sign hidden here, namely, the equality is defined by
the composition
$$
A[-i - j] \to (A \otimes_A^\mathbf{L} A)[-i - j] \to
A[-i] \otimes_A^\mathbf{L} A[-j]
$$
where in the second step we use the identification of
More on Algebra, Item (\ref{more-algebra-item-shift-tensor})
which uses a sign in principle.
Except, in this case the sign is $+1$ by our convention and even if it wasn't
$+1$ it wouldn't matter since we used the same sign
in the identification
$\mathcal{O}_X[-i - j] =
\mathcal{O}_X[-i] \otimes_{\mathcal{O}_X}^\mathbf{L} \mathcal{O}_X[-j]$.}
$$
A[-i - j] = A[-i] \otimes_A^\mathbf{L} A[-j]
\xrightarrow{\xi \otimes \eta}
R\Gamma(X, K) \otimes_A^\mathbf{L} R\Gamma(X, M)
$$
By definition the cup product $\xi \cup \eta$ is the map
$A[-i - j] \to R\Gamma(X, K \otimes_{\mathcal{O}_X}^\mathbf{L} M)$
which is adjoint to
$$
(\epsilon_K \otimes \epsilon_M) \circ Lf^*(\xi \otimes \eta) =
(\epsilon_K \circ Lf^*\xi) \otimes (\epsilon_M \circ Lf^*\eta)
$$
However, it is easy to see that
$\epsilon_K \circ Lf^*\xi = \tilde \xi$ and
$\epsilon_M \circ Lf^*\eta = \tilde \eta$.
We conclude that $\widetilde{\xi \cup \eta} = \tilde \xi \otimes \tilde \eta$
which means we have the desired agreement.
\end{proof}

\end{document}
